
Model repositories host millions of neural network checkpoints with varying architectures and tasks. Metadata is often unreliable: prior informal surveys suggest that 30-40\% of checkpoints have corrupted task labels, missing training details, or incorrect architecture tags. This metadata unreliability hinders model discovery and reuse. The central question is whether neural network weights encode task identity across diverse architectures.

Recent advances demonstrate that model properties can be extracted from parameter tensors. Task arithmetic \citep{Ilharco2022} shows that fine-tuning directions encode task-specific information. Neural Functional Networks (NFN; Zhou et al., 2023) and Universal Neural Functionals (UNF; Zhou et al., 2024) establish permutation equivariant architectures for weight space learning. However, these methods are limited to homogeneous collections. NFN operates on identical architectures, UNF handles single architecture families, and task arithmetic requires shared pretrained base checkpoints. Real model repositories are heterogeneous, containing CNNs, Transformers, RNNs, and MLPs trained from scratch on overlapping tasks.

This limitation reflects a tension between equivariance and expressivity. Neuron permutation equivariance preserves local symmetries within architectures, enabling sample-efficient learning. Extending equivariance across architectures (e.g., aligning CNN layers with Transformer blocks) requires hand-designed correspondences or abandoning equivariance entirely.

This work investigates whether task-level functional constraints create architecture-invariant structural features in weight distributions. We introduce a Hierarchical Variational Autoencoder that resolves the equivariance-expressivity tradeoff through architectural decomposition: (1) architecture-specific equivariant encoders preserve local neuron permutation symmetries, (2) permutation-invariant pooling collapses neuron-level details to layer-summary statistics, and (3) Transformer sequence modeling potentially discovers relational structure across architectures.

For proof-of-concept validation, we trained this hierarchical VAE on synthetic heterogeneous model zoo data containing 2,120 models spanning 2 architecture families (CNNs, ResNets with 4 depth variants) and 9 vision tasks. Real dataset validation is required before publication.

The hypothesis is that task constraints dominate computational primitive variance at layer-summary scale: same-task different-architecture models should cluster more tightly than different-task same-architecture models.

\textbf{Proof-of-Concept Results on Synthetic Data:} Same-task different-architecture models cluster 49.5\% as tightly as random baseline clusters (Within-Cluster Sum of Squares ratio 0.495), with strong statistical significance (p<0.000001) and large effect size (Cohen's d=1.45). Centered Kernel Alignment similarity between architecture-specific encoders exceeds 0.82 for same-task pairs (vs 0.14 for different-task pairs). Hierarchical pooling retains 68\% task-relevant information (measured via reconstruction task prediction accuracy).

These synthetic data results suggest that task-level functional constraints may create architecture-invariant structural features in weight distributions, enabling cross-architecture model property inference without shared base models or hand-designed alignment. Real ModelZooDataset validation is required to confirm external validity and rule out synthetic data artifacts.

\subsubsection{1.1 Contributions}

\begin{enumerate}
\item \textbf{Proof-of-concept demonstration on synthetic data:} Same-task different-architecture models cluster with large effect size (Cohen's d=1.45, p<0.000001) on synthetic model zoo data.
\end{enumerate}

\begin{enumerate}
\item \textbf{Hierarchical VAE design:} Three-level architecture combining architecture-specific equivariant encoders, permutation-invariant pooling, and Transformer sequence modeling.
\end{enumerate}

\begin{enumerate}
\item \textbf{Empirical validation on synthetic model zoos:} 2,120 models across 2 architecture families (4 depth variants) and 9 vision tasks, with rigorous statistical testing (bootstrap hypothesis tests, effect size analysis, feasibility gates).
\end{enumerate}

\begin{enumerate}
\item \textbf{Mechanism validation on synthetic data:} CKA feasibility gate confirms architecture subspace compatibility (0.82 same-task), reconstruction task accuracy demonstrates information preservation (68\%).
\end{enumerate}

Real dataset validation on ModelZooDataset is required before publication.

